%!TEX root = ../../main.tex

\chapter{Grundlagen}
\label{ch:grundlagen}

Dieses Kapitel legt die theoretischen Grundlagen, auf denen die experimentelle Evaluierung der folgenden Kapitel aufbaut. Zunächst werden ETL-Prozesse als etabliertes Verfahren der Datenintegration beschrieben (Abschnitt~\ref{sec:etl}). Anschließend wird das Konzept der Datenqualität als zentraler Bezugspunkt für die untersuchten ETL-Teilaufgaben eingeführt (Abschnitt~\ref{sec:datenqualitaet}). Darauf folgen eine Einordnung von LLMs als verarbeitende Komponente (Abschnitt~\ref{sec:llm}) sowie eine Darstellung des Prompt Engineering als zentrale Steuerungsmethode für diese Modelle (Abschnitt~\ref{sec:prompting}).

\section{ETL-Prozesse}
\label{sec:etl}

\subsection{Definition und Phasen}
\label{subsec:etl-phasen}

Der Begriff ETL steht für die drei aufeinanderfolgenden Phasen \emph{Extract}, \emph{Transform} und \emph{Load} und bezeichnet einen der zentralen Prozesse der Datenintegration, mit dem Daten aus unterschiedlichen Quellsystemen in ein einheitliches Zielsystem überführt werden \autocite{khanOverviewETLTechniques2024,walhaDataIntegrationTraditional2024}. Historisch ist der ETL-Prozess eng mit dem Aufbau von Data Warehouses verbunden, in denen konsolidierte und bereinigte Daten als Grundlage für Auswertungen und Entscheidungsunterstützung bereitgestellt werden \autocite{khanOverviewETLTechniques2024,boulahiaMulticriteriaEvaluationResearch2024}. Der Prozess bildet damit das Bindeglied zwischen unverarbeiteten Rohdaten und den daraus abgeleiteten, nutzbaren Informationen \autocite{khanOverviewETLTechniques2024,walhaDataIntegrationTraditional2024}.

In der ersten Phase, der Extraktion, werden die relevanten Daten aus einer oder mehreren Quellen ausgelesen. Diese Quellen sind in der Praxis heterogen und reichen von relationalen Datenbanken über Dateien und Tabellen bis hin zu Schnittstellen externer Dienste, sodass bereits hier unterschiedliche Formate und Strukturen zusammengeführt werden müssen \autocite{khanOverviewETLTechniques2024,walhaDataIntegrationTraditional2024}. Die Transformationsphase bildet den inhaltlichen Kern des Prozesses: Die extrahierten Daten werden bereinigt, vereinheitlicht und in die vom Zielsystem geforderte Struktur überführt. Typische Operationen sind die Korrektur fehlerhafter oder fehlender Werte, die Vereinheitlichung von Formaten, die Zusammenführung und Aggregation von Datensätzen, das Auflösen von Duplikaten sowie das Abbilden der Quellstrukturen auf das Zielschema \autocite{khanOverviewETLTechniques2024,walhaDataIntegrationTraditional2024,raviETLExtractTransform2025}. In der abschließenden Ladephase werden die transformierten Daten schließlich in das Zielsystem geschrieben, etwa in ein Data Warehouse oder eine Analyseplattform, wobei zwischen vollständigem Laden und inkrementellem Aktualisieren unterschieden wird \autocite{khanOverviewETLTechniques2024,boulahiaMulticriteriaEvaluationResearch2024}.

Die strikte Abfolge der drei Phasen ist dabei nicht das einzig mögliche Vorgestaltungsmuster. Insbesondere im Kontext von Big Data hat sich mit ELT (Extract, Load, Transform) eine Variante etabliert, bei der die Rohdaten zunächst geladen und erst innerhalb des Zielsystems transformiert werden, um von dessen Verarbeitungskapazitäten zu profitieren \autocite{walhaDataIntegrationTraditional2024,raviETLExtractTransform2025}. Für die vorliegende Arbeit bleibt jedoch die klassische ETL-Reihenfolge maßgeblich, da der Schwerpunkt auf den Transformationsschritten der Datenbereinigung, Datentransformation und Duplikaterkennung liegt.

\subsection{Typische Herausforderungen}
\label{subsec:etl-herausforderungen}

Obwohl ETL-Prozesse seit Jahrzehnten etabliert sind, gelten sie in modernen Datenlandschaften zunehmend als anspruchsvoll und wartungsintensiv. Diese Einschätzung gründet auf mehreren miteinander verbundenen Herausforderungen. Eine zentrale davon ist die Sicherung der Datenqualität: Da die Eingangsdaten aus unterschiedlichen Quellen stammen, weisen sie häufig Fehler, Lücken und Inkonsistenzen auf, die im Transformationsschritt erkannt und behoben werden müssen, um verlässliche Analyseergebnisse zu gewährleisten \autocite{khanOverviewETLTechniques2024,rangineniReviewEnhancingData2023}. Eng damit verbunden ist die Heterogenität der Quellen, die sich in unterschiedlichen Formaten, Schemata und Bedeutungen gleichartiger Felder äußert und einen erheblichen Aufwand bei der Vereinheitlichung verursacht \autocite{walhaDataIntegrationTraditional2024,khanOverviewETLTechniques2024}.

Hinzu kommt die Frage der Skalierbarkeit. Mit dem stetigen Wachstum von Datenvolumen und -vielfalt geraten klassische, batch-orientierte ETL-Architekturen an ihre Grenzen und können den Anforderungen an Verarbeitungsgeschwindigkeit und Aktualität nur noch eingeschränkt gerecht werden \autocite{walhaDataIntegrationTraditional2024,raviETLExtractTransform2025}. Ein weiterer kritischer Aspekt ist der hohe manuelle Aufwand bei Entwurf, Anpassung und Wartung: Klassische Pipelines sind in der Regel regelbasiert und statisch modelliert, sodass jede Änderung an den Quellstrukturen oder den fachlichen Anforderungen einen manuellen Eingriff erfordert \autocite{mishraAutomatingDataIntegration2020,nagabhyruBridgingTraditionalETL2022}. Diese geringe Anpassungsfähigkeit führt dazu, dass ein Großteil der Arbeitszeit von Datenfachleuten in die Erstellung und Pflege solcher Pipelines fließt, was sie zu einem häufig benannten Engpass datengetriebener Organisationen macht \autocite{mishraAutomatingDataIntegration2020,raviETLExtractTransform2025}. Genau diese Schwachstellen motivieren die in dieser Arbeit untersuchte Frage, inwieweit lernende und sprachbasierte Verfahren den manuellen Aufwand reduzieren und die Flexibilität klassischer ETL-Prozesse erhöhen können.

\subsection{Klassische Werkzeuge und Verfahren}
\label{subsec:etl-werkzeuge}

Zur Umsetzung von ETL-Prozessen hat sich über die Jahre ein breites Spektrum an Werkzeugen und Verfahren herausgebildet. Auf der einen Seite stehen dedizierte ETL-Plattformen mit grafischer Modellierung wie Talend, Pentaho oder Informatica, die den Prozess über vorkonfigurierte Komponenten abbilden und sich vor allem in betrieblichen Data-Warehouse-Umgebungen etabliert haben \autocite{khanOverviewETLTechniques2024,boulahiaMulticriteriaEvaluationResearch2024}. Auf der anderen Seite werden ETL-Strecken zunehmend programmatisch umgesetzt, etwa mit der Python-Bibliothek Pandas für die Verarbeitung tabellarischer Daten oder mit verteilten Frameworks wie Apache Spark für große Datenmengen \autocite{khanOverviewETLTechniques2024,walhaDataIntegrationTraditional2024}.

Allen klassischen Ansätzen ist gemeinsam, dass die Verarbeitungslogik regelbasiert und explizit vorgegeben wird: Entwicklerinnen und Entwickler definieren die Transformations- und Bereinigungsregeln manuell, etwa in Form von Bedingungen, Zuordnungstabellen oder Skripten \autocite{mishraAutomatingDataIntegration2020,boulahiaMulticriteriaEvaluationResearch2024}. Dieser Ansatz bietet den Vorteil hoher Nachvollziehbarkeit und deterministischer, reproduzierbarer Ergebnisse, stößt jedoch dort an Grenzen, wo Datenprobleme vielfältig, kontextabhängig oder schwer im Voraus zu formalisieren sind \autocite{mishraAutomatingDataIntegration2020,nagabhyruBridgingTraditionalETL2022}. In der vorliegenden Arbeit dient eine solche regelbasiert implementierte Pipeline als Vergleichsmaßstab, gegen den die untersuchten LLM-gestützten Verfahren gemessen werden.

\section{Datenqualität und Datenbereinigung}
\label{sec:datenqualitaet}

\subsection{Dimensionen der Datenqualität}
\label{subsec:dq-dimensionen}

Datenqualität ist ein mehrdimensionales Konstrukt, das beschreibt, inwieweit ein Datenbestand für seinen vorgesehenen Verwendungszweck geeignet ist. In der Literatur hat sich ein Kern wiederkehrender Qualitätsdimensionen herausgebildet, zu denen insbesondere Genauigkeit, Vollständigkeit, Konsistenz, Aktualität und Eindeutigkeit zählen \autocite{ridzuanReviewDataQuality2024,rangineniReviewEnhancingData2023}. Die Genauigkeit beschreibt, inwieweit die gespeicherten Werte die reale Sachlage korrekt abbilden, während die Vollständigkeit den Grad angibt, in dem alle erforderlichen Werte tatsächlich vorhanden sind \autocite{ridzuanReviewDataQuality2024,rangineniReviewEnhancingData2023}. Die Konsistenz bezieht sich auf die Widerspruchsfreiheit der Daten innerhalb eines Bestands und über Systemgrenzen hinweg, etwa bei einheitlichen Formaten und Bezeichnungen \autocite{ridzuanReviewDataQuality2024,rangineniReviewEnhancingData2023}. Die Aktualität erfasst, ob die Daten den jeweils gültigen Zeitstand widerspiegeln, und die Eindeutigkeit verlangt, dass jeder reale Sachverhalt nur einmal und nicht in mehreren widersprüchlichen oder duplizierten Datensätzen abgelegt ist \autocite{ridzuanReviewDataQuality2024,rangineniReviewEnhancingData2023}.

Mit dem Aufkommen von Big Data ist dieser klassische Kanon erweitert worden, da Volumen, Geschwindigkeit und Vielfalt der Daten zusätzliche Anforderungen stellen und feinere Differenzierungen der Qualitätsdimensionen erforderlich machen \autocite{ridzuanReviewDataQuality2024,rangineniReviewEnhancingData2023}. Für die vorliegende Arbeit sind insbesondere die Dimensionen Genauigkeit, Vollständigkeit, Konsistenz und Eindeutigkeit von Bedeutung, da sie unmittelbar mit den untersuchten ETL-Teilaufgaben der Datenbereinigung, Datentransformation und Duplikaterkennung korrespondieren.

\subsection{Häufige Datenprobleme}
\label{subsec:dq-probleme}

Aus den Qualitätsdimensionen lassen sich die in der Praxis am häufigsten auftretenden Datenprobleme ableiten, die im Rahmen der Datenbereinigung erkannt und behoben werden müssen. Fehlende Werte gehören zu den verbreitetsten Mängeln und beeinträchtigen unmittelbar die Vollständigkeit; sie erfordern Entscheidungen über Imputation, Entfernung oder Kennzeichnung der betroffenen Datensätze \autocite{coteDataCleaningMachine2024,rangineniReviewEnhancingData2023}. Inkonsistente Formate – etwa unterschiedliche Datums-, Zahlen- oder Einheitendarstellungen sowie abweichende Schreibweisen gleicher Inhalte – beeinträchtigen die Konsistenz und müssen im Transformationsschritt vereinheitlicht werden \autocite{rangineniReviewEnhancingData2023,ridzuanReviewDataQuality2024}.

Ein weiteres zentrales Problem sind Duplikate, also mehrfach vorhandene Datensätze, die denselben realen Sachverhalt beschreiben, sich aber in Schreibweise oder Vollständigkeit unterscheiden können. Ihre Erkennung ist anspruchsvoll, da sie über den exakten Abgleich hinaus häufig eine unscharfe Ähnlichkeitsbetrachtung erfordert \autocite{coteDataCleaningMachine2024,maddaliAutomatingDataQuality}. Hinzu kommen Ausreißer und fehlerhafte Werte, die durch Eingabefehler, Messfehler oder Systembrüche entstehen und die Genauigkeit des Datenbestands gefährden \autocite{coteDataCleaningMachine2024,rangineniReviewEnhancingData2023}. Die Bedeutung einer gründlichen Bereinigung ergibt sich daraus, dass mangelhafte Eingangsdaten unmittelbar auf nachgelagerte Analysen und Modelle durchschlagen, was häufig mit dem Grundsatz \enquote{garbage in, garbage out} umschrieben wird \autocite{ridzuanReviewDataQuality2024,rangineniReviewEnhancingData2023}. Die gezielte Konstruktion solcher Datenprobleme bildet später die Grundlage des in Kapitel~5 beschriebenen synthetischen Datensatzes.

\section{Large Language Models}
\label{sec:llm}

\subsection{Funktionsweise und Architektur}
\label{subsec:llm-architektur}

LLMs sind Sprachmodelle mit einer sehr großen Zahl an Parametern, die auf großen Mengen an Textdaten vortrainiert werden und in der Lage sind, natürliche Sprache zu verstehen und zu erzeugen \autocite{zhaoSurveyLargeLanguage2026,liangSurveyMultimodelLarge2024}. Technische Grundlage nahezu aller aktuellen Modelle ist die Transformer-Architektur, deren Selbstaufmerksamkeitsmechanismus es erlaubt, Abhängigkeiten zwischen weit auseinanderliegenden Textbestandteilen effizient zu erfassen \autocite{zhaoSurveyLargeLanguage2026,liangSurveyMultimodelLarge2024}. Eingaben werden dabei zunächst in kleinere Einheiten, sogenannte Token, zerlegt, die das Modell intern als numerische Vektoren verarbeitet; die Ausgabe entsteht, indem das Modell schrittweise das jeweils wahrscheinlichste nachfolgende Token vorhersagt \autocite{zhaoSurveyLargeLanguage2026}.

Charakteristisch für LLMs ist der zweistufige Entstehungsprozess aus einem aufwendigen Vortraining auf großen, allgemeinen Datenmengen und einer anschließenden Anpassung, etwa durch Feinabstimmung oder durch ein Training mit menschlichem Feedback, das die Modelle an die Erwartungen der Nutzenden ausrichtet \autocite{zhaoSurveyLargeLanguage2026}. Ein zentraler Befund der Forschung ist, dass mit zunehmender Modell- und Datengröße neue Fähigkeiten auftreten, die kleinere Modelle nicht aufweisen. Dazu zählt insbesondere das In-Context-Lernen, also die Fähigkeit, eine Aufgabe allein anhand von Anweisungen und Beispielen im Eingabetext zu lösen, ohne dass die Modellparameter verändert werden \autocite{zhaoSurveyLargeLanguage2026,liuPretrainPromptPredict2023}. Genau diese Eigenschaft bildet die Grundlage dafür, LLMs ohne aufgabenspezifisches Nachtraining für ETL-Teilaufgaben einzusetzen. Zugleich ist die Ausgabe der Modelle probabilistischer Natur: Da das nächste Token aus einer Wahrscheinlichkeitsverteilung gezogen wird, können identische Eingaben zu unterschiedlichen Ausgaben führen, was unmittelbare Folgen für die Reproduzierbarkeit der Ergebnisse hat \autocite{zhaoSurveyLargeLanguage2026}.

\subsection{Überblick aktueller Modelle}
\label{subsec:llm-modelle}

Seit der breiten öffentlichen Aufmerksamkeit für dialogfähige Modelle hat sich das Feld der LLMs rasant entwickelt und ausdifferenziert \autocite{zhaoSurveyLargeLanguage2026,2026AIIndex}. Auf dem Markt konkurriert heute eine Reihe leistungsfähiger Modellfamilien, darunter die GPT-Modelle von OpenAI, die Claude-Modelle von Anthropic, die Gemini-Modelle von Google sowie offen verfügbare Modelle wie die Llama-Reihe \autocite{zhaoSurveyLargeLanguage2026,2026AIIndex}. Diese Modelle unterscheiden sich nicht nur in ihrer Architektur und Parameterzahl, sondern auch in praxisrelevanten Eigenschaften wie der Größe des Kontextfensters, der Verarbeitungsgeschwindigkeit, den Nutzungskosten und der Verfügbarkeit als offenes oder geschlossenes Modell \autocite{zhaoSurveyLargeLanguage2026,2026AIIndex}.

Eine wesentliche Entwicklung der jüngeren Vergangenheit ist die Erweiterung hin zu multimodalen Modellen, die neben Text auch weitere Modalitäten wie Bilder verarbeiten können \autocite{liangSurveyMultimodelLarge2024,zhaoSurveyLargeLanguage2026}. Angesichts der hohen Innovationsgeschwindigkeit verschieben sich die Leistungsgrenzen zwischen den Modellen kontinuierlich, weshalb vergleichende Bewertungen anhand standardisierter Kennzahlen wie Qualität, Geschwindigkeit und Preis an Bedeutung gewonnen haben \autocite{ComparisonAIModels,2026AIIndex}. Für die vorliegende Arbeit ist diese Heterogenität von unmittelbarer Relevanz, da sie die Auswahl mehrerer aktueller Modelle für den experimentellen Vergleich begründet und nahelegt, dass sich die Eignung für ETL-Aufgaben zwischen den Modellen unterscheiden kann.

\subsection{Stärken und Limitationen bei strukturierten Daten}
\label{subsec:llm-strukturierte-daten}

Im Kontext der Datenverarbeitung liegt die Stärke von LLMs in ihrer Flexibilität: Sie können semi- und unstrukturierte Inhalte interpretieren, kontextabhängige Zusammenhänge erfassen und über reine Mustererkennung hinaus auch implizites Wissen einbringen, ohne dass die zugrunde liegenden Regeln explizit vorgegeben werden müssen \autocite{hassaniRoleChatGPTData2023,zhaoSurveyLargeLanguage2026}. Damit eignen sie sich potenziell für Aufgaben wie die Bereinigung uneinheitlicher Werte, die Vereinheitlichung von Formaten oder das Erkennen ähnlicher Datensätze, die mit starren Regeln nur schwer abzudecken sind \autocite{hassaniRoleChatGPTData2023}.

Diesen Stärken stehen jedoch erhebliche Limitationen gegenüber. Eine zentrale Einschränkung ist die Neigung zu Halluzinationen, also zur Erzeugung plausibel wirkender, aber sachlich falscher oder frei erfundener Ausgaben, was bei der Verarbeitung strukturierter Daten zu schwer erkennbaren Fehlern führen kann \autocite{hassaniRoleChatGPTData2023,zhaoSurveyLargeLanguage2026}. Hinzu kommt die Begrenzung durch das Kontextfenster, die der Menge an Daten, die in einem einzelnen Verarbeitungsschritt übergeben werden kann, eine technische Obergrenze setzt \autocite{zhaoSurveyLargeLanguage2026}. Ferner fallen die Nutzungskosten und die Verarbeitungszeit ins Gewicht, die insbesondere bei großen Datenmengen zu einem begrenzenden Faktor werden \autocite{zhaoSurveyLargeLanguage2026,2026AIIndex}. Schließlich erschwert die bereits beschriebene probabilistische Natur der Modelle die Reproduzierbarkeit, da wiederholte Ausführungen identischer Aufgaben voneinander abweichende Ergebnisse liefern können \autocite{zhaoSurveyLargeLanguage2026,hassaniRoleChatGPTData2023}. Diese Spannung zwischen Flexibilität und Verlässlichkeit bildet einen der zentralen Ausgangspunkte und zugleich einen der Untersuchungsgegenstände der vorliegenden Arbeit.

\section{Prompt Engineering}
\label{sec:prompting}

\subsection{Grundlegende Prompting-Strategien}
\label{subsec:prompting-strategien}

Da LLMs ihre Aufgaben primär über das In-Context-Lernen lösen, kommt der Gestaltung der Eingabe – dem sogenannten Prompt – eine entscheidende Bedeutung zu. Unter Prompt Engineering versteht man das gezielte Entwerfen und Optimieren solcher Eingaben, um das Modellverhalten zu steuern, ohne die Modellparameter selbst zu verändern \autocite{sahooSystematicSurveyPrompt,schulhoffPromptReportSystematic2025}. Die einfachste Strategie ist das Zero-Shot-Prompting, bei dem das Modell eine Aufgabe allein anhand einer Anweisung und ohne weitere Beispiele bearbeitet \autocite{liuPretrainPromptPredict2023,sahooSystematicSurveyPrompt}. Beim Few-Shot-Prompting werden der Anweisung hingegen einige beispielhafte Ein- und Ausgabepaare beigefügt, anhand derer das Modell das gewünschte Verhalten ableitet; dieser Ansatz verbessert die Ergebnisqualität bei vielen Aufgaben spürbar \autocite{liuPretrainPromptPredict2023,sahooSystematicSurveyPrompt}.

Darüber hinaus haben sich anspruchsvollere Strategien etabliert. Beim Chain-of-Thought-Prompting wird das Modell angeleitet, seine Lösung in nachvollziehbaren Zwischenschritten zu entwickeln, was sich insbesondere bei mehrstufigen oder schlussfolgernden Aufgaben als wirksam erwiesen hat \autocite{sahooSystematicSurveyPrompt,schulhoffPromptReportSystematic2025}. Eine weitere verbreitete Technik ist das Role-Prompting, bei dem dem Modell eine bestimmte Rolle oder Perspektive zugewiesen wird, um Stil und Inhalt der Ausgabe zu lenken \autocite{schulhoffPromptReportSystematic2025,sahooSystematicSurveyPrompt}. Angesichts der Vielzahl mittlerweile beschriebener Techniken weisen systematische Übersichten allerdings auf eine uneinheitliche Begriffsbildung und eine teils fragmentierte Forschungslage hin, was eine sorgfältige Auswahl und Definition der eingesetzten Strategien erforderlich macht \autocite{schulhoffPromptReportSystematic2025,sahooSystematicSurveyPrompt}. Für die vorliegende Arbeit werden daher klar abgegrenzte Prompting-Strategien definiert und systematisch gegeneinander verglichen.

\subsection{Einflussfaktoren auf die Ergebnisqualität}
\label{subsec:prompting-einfluss}

Die Qualität der von einem LLM erzeugten Ausgabe hängt von einer Reihe von Einflussfaktoren ab, die über die reine Wahl der Prompting-Strategie hinausgehen. Beim Few-Shot-Prompting wirken sich Auswahl, Anzahl und Reihenfolge der mitgegebenen Beispiele unmittelbar auf das Ergebnis aus, sodass schon kleine Änderungen an den Beispielen die Modellantwort verschieben können \autocite{schulhoffPromptReportSystematic2025,liuPretrainPromptPredict2023}. Auch die genaue Formulierung und Strukturierung der Anweisung beeinflusst die Ergebnisqualität erheblich, weshalb das Auffinden geeigneter Prompts häufig einen iterativen Prozess darstellt \autocite{sahooSystematicSurveyPrompt,changEfficientPromptingMethods2024}.

Ein technisch bedeutsamer Parameter ist die Temperatur, die steuert, wie stark die Ausgabe von der wahrscheinlichsten Vorhersage abweicht: Niedrige Werte führen zu deterministischeren und stabileren, hohe Werte zu vielfältigeren, aber weniger vorhersagbaren Antworten, was die Reproduzierbarkeit der Ergebnisse unmittelbar berührt \autocite{schulhoffPromptReportSystematic2025,zhaoSurveyLargeLanguage2026}. Zudem skaliert die Leistungsfähigkeit beim Prompting mit der Größe und Leistungsfähigkeit des zugrunde liegenden Modells, sodass dieselbe Strategie je nach Modell unterschiedlich gut wirkt \autocite{liuPretrainPromptPredict2023,schulhoffPromptReportSystematic2025}. Schließlich ist mit längeren und beispielreicheren Prompts ein höherer Verbrauch an Token und damit ein höherer Aufwand verbunden, sodass zwischen Ergebnisqualität und Effizienz abzuwägen ist \autocite{changEfficientPromptingMethods2024,schulhoffPromptReportSystematic2025}. Diese Faktoren begründen, warum in der vorliegenden Arbeit sowohl unterschiedliche Modelle als auch unterschiedliche Prompting-Strategien unter kontrollierten Bedingungen verglichen und ihre Ergebnisse mehrfach erhoben werden.
